\documentclass[a4paper]{article}

%% Language and font encodings
\usepackage[english]{babel}
\usepackage[utf8]{inputenc}
\usepackage[T1]{fontenc}

%% Page size and margins
\usepackage[a4paper,top=2cm,bottom=2cm,left=2.5cm,right=2.5cm,marginparwidth=1.75cm]{geometry}

%% Core packages
\usepackage{amsmath}
\usepackage{graphicx}
\usepackage{tabularx}
\usepackage{booktabs}
\usepackage{float}
\usepackage{subcaption}
\usepackage{microtype}
\usepackage{xcolor}
\definecolor{tcdBlue}{RGB}{5, 105, 185}

%% Hyperlinks
\usepackage[colorlinks=true,allcolors=.,urlcolor=blue]{hyperref}

%% Section heading style
\renewcommand{\familydefault}{\sfdefault}
\usepackage{titlesec}
\titleformat{\section}[hang]{\normalfont\Large\bfseries\color{tcdBlue}}{\thesection}{1em}{}{}

\hyphenation{op-tical net-works semi-conduc-tor}

\begin{document}

\title{\textbf{Representation Learning in Deep Reinforcement Learning:}\\
       \large A Comparative Study of DQN and Double DQN on Atari}
\author{Vinol Tauro\\
        \small Student Number: 25340566\\
        \small School of Computer Science and Statistics\\
        \small Trinity College Dublin}
\date{June 2026}
\maketitle

\begin{abstract}
This study investigates how structurally similar yet strategically distinct games shape the internal representations of deep reinforcement learning agents, and whether the choice of training algorithm affects what agents learn to represent. Pong and Breakout share fundamental low-level visual structure, including a ball, a paddle, and deflection physics, but differ substantially in high-level strategy. Using a 2$\times$2 factorial design, Deep Q-Networks (DQN) and Double DQN (DDQN) agents were trained on both games and their 512-dimensional penultimate layer activations were analysed using t-SNE, Grad-CAM saliency maps, dead neuron analysis, cross-game cosine similarity, and layer-wise Centred Kernel Alignment (CKA). The results show that despite shared visual structure, game-specific strategic demands drive representational divergence: agents playing different games develop clearly distinct internal encodings. CKA analysis confirms a representational hierarchy -- early convolutional layers remain highly similar across games while the representation layer diverges substantially -- providing quantitative evidence that generalisation and specialisation are spatially separated within the network. Non-trivial cross-game cosine similarity persists, reflecting shared low-level visual features, and this alignment is stronger under DDQN. Within each game, DDQN produces more compact and structured representations, consistent with its cleaner gradient signal reducing representational noise. These findings reveal how the boundary between shared and game-specific representations aligns with the boundary between visual and strategic game features, and how that boundary can be localised to specific network layers.
\end{abstract}

\tableofcontents
\newpage

% ============================================================
\section{Introduction}
% ============================================================

Pong and Breakout look similar on the surface. Both feature a ball bouncing off a paddle, both share the same deflection physics, and both are rendered in the same pixelated Atari style. Yet anyone who has played both knows they demand completely different things from a player: Pong is about reading and reacting to an opponent, Breakout is about planning a path through a grid of bricks. The question this study asks is whether a deep reinforcement learning agent, trained purely from pixels and reward signals, picks up on that distinction internally, and if so, where the line falls between what is shared and what is game-specific in its learned representations.

Understanding the internal representations of deep RL agents matters beyond academic interest. As these systems are applied in more complex and consequential settings, interpretability becomes a practical concern. Knowing whether an agent has built a clean, structured model of its environment, or a noisy approximation that happens to score well, has direct implications for how much we should trust it and how well it will transfer to new tasks.

This study trains DQN \cite{mnih2015human} and Double DQN (DDQN) \cite{van2016deep} agents on both games and examines the 512-dimensional activations at the penultimate layer of the shared CNN backbone. DQN and DDQN are used not because the algorithm comparison is the primary question, but because having two independent algorithms provides two separate windows into the same representational phenomenon. If a finding holds across both, it is more likely to reflect something genuine about the games rather than an artefact of a particular training procedure.

Pong and Breakout were chosen specifically because they sit at an interesting point on the similarity spectrum: visually similar enough that one might expect substantial representational overlap, yet strategically different enough that divergence would not be surprising. This makes them a useful pair for probing exactly where shared structure ends and game-specific structure begins.

Three questions guide the work:
\begin{enumerate}
    \item Do agents trained on Pong and Breakout develop similar internal representations, given the shared low-level visual structure of the two games?
    \item Where does representational similarity break down, and does this correspond to the boundary between shared visual features and game-specific strategic demands?
    \item Does the training algorithm, DQN versus DDQN, affect the quality and structure of the learned representations within each game?
\end{enumerate}

% ============================================================
\section{Background}
\label{sec:background}
% ============================================================

\subsection{Deep Q-Networks}

DQN \cite{mnih2015human} approximates the action-value function $Q(s, a)$, the expected cumulative reward for taking action $a$ in state $s$, using a convolutional neural network. Two mechanisms make this tractable for raw pixel inputs.

The first is \textit{experience replay}: transitions $(s, a, r, s')$ are stored in a circular buffer and sampled randomly during training, breaking the temporal correlations between consecutive frames that would otherwise destabilise learning.

The second is a \textit{target network}: a periodically frozen copy of the online network used to compute training targets. Without it, targets shift with every gradient update, making training unstable. The DQN TD target is:

\begin{equation}
    y = r + \gamma \cdot \max_{a'} Q(s', a';\, \theta^{-})
\end{equation}

where $\theta^{-}$ are the target network parameters. Critically, the same network $\theta^{-}$ both selects the best next action and evaluates its value. Taking the maximum over noisy Q-value estimates produces a systematic upward bias that compounds over training.

\subsection{Double DQN}

Van Hasselt et al.\ \cite{van2016deep} proposed a straightforward fix: use the online network to select the best next action and the target network to evaluate it. The DDQN target is:

\begin{equation}
    y = r + \gamma \cdot Q\!\left(s',\, \argmax_{a'} Q(s', a';\, \theta);\, \theta^{-}\right)
\end{equation}

Because the two networks have different error profiles at any point during training, their errors are unlikely to compound in the same direction, substantially reducing overestimation bias. In this study's codebase, \texttt{DoubleDQNAgent} inherits everything from \texttt{DQNAgent} and overrides only the \texttt{learn()} method. Any representational difference observed between the two agents is attributable solely to this target formulation.

\subsection{Representation Learning and Capacity Loss}

The 512-dimensional activation vector at the penultimate fully connected layer serves as the agent's compressed encoding of the current game state. Zahavy et al.\ \cite{zahavy2016graying} applied t-SNE to DQN representations and showed that semantically similar game states cluster together in the embedding space. This study extends that approach to a controlled multi-algorithm, multi-game comparison.

Lyle et al.\ \cite{lyle2022understanding} demonstrated that deep RL agents frequently lose representational capacity during training: neurons become permanently inactive due to the ReLU dead zone, reducing the effective dimensionality of the representation layer. Noisier gradient signals, of the kind produced by DQN's overestimation bias, accelerate this process. This provides a direct theoretical motivation for expecting DDQN to maintain a more active and expressive representation layer.

\subsection{t-SNE}

t-SNE \cite{van2008visualizing} reduces high-dimensional data to two dimensions by minimising the KL divergence between Gaussian neighbourhood distributions in the original space and Student-t distributions in the 2D projection. It faithfully represents local clustering, though inter-cluster distances in the projection are not proportional to distances in the original high-dimensional space. Qualitative cluster structure is therefore interpreted visually from t-SNE plots, while quantitative assessments use the silhouette score and Davies-Bouldin index computed directly in the 512-dimensional space.

\subsection{Grad-CAM}

Grad-CAM \cite{selvaraju2017grad} produces visual explanations of a network's decisions by backpropagating gradients to the final convolutional layer and weighting feature maps by their importance to the output. Applied to Atari agents, it highlights which regions of the game frame most influence action selection at any given step.

\subsection{Centred Kernel Alignment}

Centred Kernel Alignment (CKA) \cite{kornblith2019similarity} measures the similarity between two sets of representations, independent of linear transformations. Given activation matrices $X \in \mathbb{R}^{n \times p}$ and $Y \in \mathbb{R}^{n \times q}$ collected over the same $n$ inputs, linear CKA is defined as:

\begin{equation}
    \text{CKA}(X, Y) = \frac{\|Y^{\top}X\|_F^2}{\|X^{\top}X\|_F \cdot \|Y^{\top}Y\|_F}
\end{equation}

CKA returns values in $[0, 1]$, where 1 indicates identical representational geometry up to linear transformation and 0 indicates complete dissimilarity. Unlike cosine similarity between mean vectors, CKA captures the full distributional geometry of the representation space at each layer, making it well-suited for comparing how similar two networks are at different depths.

% ============================================================
\section{Experimental Design}
\label{sec:method}
% ============================================================

\subsection{Factorial Structure}

The study uses a fully crossed 2$\times$2 factorial design, crossing algorithm (DQN, DDQN) with game (Pong, Breakout):

\begin{table}[H]
\centering
\caption{Factorial design. All four runs use identical hyperparameters and the same random seed.}
\label{tab:factorial}
\begin{tabular}{lcc}
\toprule
 & \textbf{Pong} & \textbf{Breakout} \\
\midrule
\textbf{DQN}  & Run 1 (2M steps) & Run 2 (5M steps) \\
\textbf{DDQN} & Run 3 (2M steps) & Run 4 (5M steps) \\
\bottomrule
\end{tabular}
\end{table}

Pong converges relatively quickly and was trained for 2M steps; Breakout is more complex and required 5M. The design isolates three comparisons: the game effect (same algorithm, different game), the algorithm effect (same game, different algorithm), and their interaction.

\subsection{Preprocessing Pipeline}

Raw 210$\times$160 RGB frames are transformed through the standard DeepMind preprocessing stack \cite{mnih2015human}:
\begin{itemize}
    \item \textbf{No-op reset:} 1 to 30 random no-op actions at the start of each episode, preventing the agent from memorising fixed starting states.
    \item \textbf{Frame skip:} the same action is repeated for 4 frames, with pixel-wise max-pooling over the last two to handle sprite flickering.
    \item \textbf{Episodic life:} each life loss is treated as a terminal signal, providing a denser learning signal.
    \item \textbf{Grayscale and resize:} frames are converted to 84$\times$84 grayscale.
    \item \textbf{Reward clipping:} rewards are clipped to $\{-1, 0, +1\}$, normalising scale across games.
    \item \textbf{Frame stacking:} the last 4 frames are stacked along the channel dimension, providing the agent with implicit velocity information.
\end{itemize}

\subsection{Network Architecture}

Both agents share the CNN from Mnih et al.\ \cite{mnih2015human}: three convolutional layers (32 filters, 8$\times$8, stride 4; 64 filters, 4$\times$4, stride 2; 64 filters, 3$\times$3, stride 1), followed by a 512-dimensional fully connected representation layer with ReLU activation, and a linear output head producing Q-values for each action. A PyTorch forward hook on the representation layer captures the 512-dim activation after every forward pass without modifying the computation. Weights are initialised using orthogonal initialisation \cite{saxe2013exact}, which produces more stable early learning than the default Xavier scheme.

\subsection{Hyperparameters}

\begin{table}[H]
\centering
\caption{Hyperparameters used across all four runs.}
\label{tab:hyperparams}
\begin{tabularx}{\linewidth}{lX}
\toprule
\textbf{Parameter} & \textbf{Value} \\
\midrule
Learning rate            & $1 \times 10^{-4}$ (Adam) \\
Replay buffer            & 100,000 transitions \\
Batch size               & 32 \\
Discount factor $\gamma$ & 0.99 \\
Target update frequency  & 1,000 steps (hard copy) \\
Epsilon schedule         & $1.0 \to 0.01$ over 100,000 steps \\
Gradient clipping        & $L_2$ norm $\leq 10.0$ \\
Frame skip / stack       & 4 / 4 \\
Input resolution         & 84$\times$84 grayscale \\
Checkpoint frequency     & Every 500,000 steps \\
Random seed              & 42 \\
\bottomrule
\end{tabularx}
\end{table}

The replay buffer of 100,000 transitions is smaller than the 1M used in the original DQN paper, a practical concession to memory constraints. All four runs use an identical buffer size, so the reduction does not confound the algorithm comparison.

\subsection{Representation Extraction}

After training, each checkpoint is loaded and the agent runs in near-greedy mode ($\varepsilon = 0.05$) for 5,000 steps. At each step the 512-dim representation vector, action taken, and cumulative episode reward are recorded, producing one \texttt{.npz} file per checkpoint and capturing the full developmental trajectory of the representation space across training.

\subsection{Layer-wise CKA Analysis}

To localise where cross-game representational similarity breaks down within the network, forward hooks are registered at four depths: the output of each of the three convolutional layers (conv1, conv2, conv3) and the 512-dimensional fully connected representation layer (fc\_repr). For a fixed set of game states, activation matrices are collected at each depth for DQN/Pong and DQN/Breakout, and pairwise CKA is computed between the two agents. The same analysis is repeated for DDQN/Pong and DDQN/Breakout. This produces a depth profile of cross-game similarity, revealing at which layer the shared visual structure gives way to game-specific strategic encoding.

% ============================================================
\section{Results}
\label{sec:results}
% ============================================================

\subsection{Training Performance}

All four agents were trained to completion. Table~\ref{tab:results} summarises the final performance and Q-value statistics for each run.

\begin{table}[H]
\centering
\caption{Training results for all four agents. Reward statistics are computed over the final 50 episodes. Peak Q is the maximum mean Q-value recorded across the full training run.}
\label{tab:results}
\begin{tabularx}{\linewidth}{lXXXX}
\toprule
\textbf{Agent} & \textbf{Episodes} & \textbf{Final Reward} & \textbf{Final Q} & \textbf{Peak Q} \\
\midrule
DQN / Pong      & 2,434   & $7.12 \pm 5.22$ & 2.64 & 2.95 \\
DDQN / Pong     & 2,457   & $7.50 \pm 5.03$ & 2.05 & 2.48 \\
DQN / Breakout  & 144,751 & $3.86 \pm 4.94$ & 4.06 & 5.39 \\
DDQN / Breakout & 144,948 & $4.00 \pm 5.04$ & 4.13 & 4.55 \\
\bottomrule
\end{tabularx}
\end{table}

On Pong, both algorithms reach similar final reward: DQN achieves $7.12 \pm 5.22$ and DDQN $7.50 \pm 5.03$, with DDQN showing a marginal advantage. The learning trajectory follows the expected pattern, starting at $-21$ (losing every point), crossing into positive reward around 1M steps, and stabilising through late training. On Breakout, DQN reaches $3.86 \pm 4.94$ and DDQN $4.00 \pm 5.04$ over 5M steps, with both algorithms reaching comparable final performance.

The near-identical task performance between DQN and DDQN is itself meaningful. Any representational differences observed in subsequent analysis cannot be attributed to differences in task performance. Both algorithms learned to play at the same level, yet as the following sections show, they did so through representations of noticeably different quality.

\begin{figure}[H]
    \centering
    \begin{subfigure}[b]{0.48\linewidth}
        \includegraphics[width=\linewidth]{figures/training_curves_pong.png}
        \caption{Pong (20-episode rolling average).}
    \end{subfigure}
    \hfill
    \begin{subfigure}[b]{0.48\linewidth}
        \includegraphics[width=\linewidth]{figures/training_curves_breakout.png}
        \caption{Breakout (20-episode rolling average).}
    \end{subfigure}
    \caption{Training curves for all four agents. Both algorithms reach comparable final reward on both games.}
    \label{fig:training_curves}
\end{figure}

\subsection{Q-Value Overestimation}

The mean maximum Q-value diverges clearly between DQN and DDQN over training (Figure~\ref{fig:qvalue}). On Pong, DQN Q-values peak at 2.95 and settle at 2.64, while DDQN peaks at 2.48 and settles at 2.05, a difference of approximately 29\% in final mean Q-value. On Breakout the gap is smaller: DQN reaches a final Q of 4.06 versus 4.13 for DDQN, with DQN showing a higher peak (5.39 versus 4.55), consistent with overestimation inflating Q-values during the most uncertain phases of early Breakout training.

The cleaner gradient signal under DDQN is consistent with the core finding of van Hasselt et al.\ \cite{van2016deep} and provides a direct mechanism for the representational differences documented in the following sections.

\begin{figure}[H]
    \centering
    \includegraphics[width=0.7\linewidth]{figures/qvalue_overestimation.png}
    \caption{Mean maximum Q-value over training for DQN and DDQN on both games. DQN values drift upward throughout training; DDQN values remain considerably lower and more stable.}
    \label{fig:qvalue}
\end{figure}

\subsection{t-SNE Representation Analysis}

\subsubsection{Game Effect}

The most prominent structural feature of the t-SNE embeddings is separation by game. When DQN/Pong and DQN/Breakout representations are projected together, two distinct clusters emerge corresponding to each game (Figure~\ref{fig:game_effect}). The same separation appears for DDQN. Regardless of which algorithm trained the agent, the game it played is the primary organising principle of its representation space.

\begin{figure}[H]
    \centering
    \begin{subfigure}[b]{0.48\linewidth}
        \includegraphics[width=\linewidth]{figures/tsne_game_effect_dqn.png}
        \caption{DQN: Pong (blue) vs Breakout (orange).}
    \end{subfigure}
    \hfill
    \begin{subfigure}[b]{0.48\linewidth}
        \includegraphics[width=\linewidth]{figures/tsne_game_effect_ddqn.png}
        \caption{DDQN: Pong vs Breakout.}
    \end{subfigure}
    \caption{Game effect: representations separate cleanly by game regardless of algorithm.}
    \label{fig:game_effect}
\end{figure}

\subsubsection{Algorithm Effect}

Within each game, the algorithm produces observable differences. On Pong, DQN and DDQN representations overlap in t-SNE space but DDQN shows noticeably tighter cluster geometry with lower intra-cluster variance, consistent with a cleaner gradient signal producing more organised internal representations (Figure~\ref{fig:algo_effect}). The same pattern appears on Breakout, though less pronounced, suggesting the algorithm effect is partially absorbed by the greater environmental complexity of that game.

\begin{figure}[H]
    \centering
    \begin{subfigure}[b]{0.48\linewidth}
        \includegraphics[width=\linewidth]{figures/tsne_algo_effect_pong.png}
        \caption{Algorithm effect on Pong.}
    \end{subfigure}
    \hfill
    \begin{subfigure}[b]{0.48\linewidth}
        \includegraphics[width=\linewidth]{figures/tsne_algo_effect_breakout.png}
        \caption{Algorithm effect on Breakout.}
    \end{subfigure}
    \caption{Algorithm effect: DDQN produces tighter cluster geometry within each game.}
    \label{fig:algo_effect}
\end{figure}

The joint projection of all four agents (Figure~\ref{fig:all_agents}) makes the hierarchy clear: two primary clusters by game, with algorithm-level sub-structure visible within each. The representation space is organised first by what the agent learned about, and second by how it learned.

\begin{figure}[H]
    \centering
    \includegraphics[width=0.6\linewidth]{figures/tsne_all_agents.png}
    \caption{All four agents projected together. Game clusters dominate; algorithm sub-structure is visible within each cluster.}
    \label{fig:all_agents}
\end{figure}

Representations coloured by cumulative episode reward (Figure~\ref{fig:reward}) show a partial spatial correspondence between t-SNE position and value, more consistently under DDQN, suggesting DDQN representations encode some value-relevant structure beyond purely perceptual features.

\begin{figure}[H]
    \centering
    \includegraphics[width=0.65\linewidth]{figures/tsne_by_reward.png}
    \caption{t-SNE coloured by cumulative reward. High-reward states (green) and low-reward states (red) show spatial correspondence in the embedding, more consistently under DDQN.}
    \label{fig:reward}
\end{figure}

\subsubsection{Temporal Evolution}

Figure~\ref{fig:temporal} shows how the representation space evolves across training checkpoints for all four agents. Early representations form a diffuse, unstructured cloud. Clusters emerge and tighten as training progresses. DDQN shows earlier emergence of clear structure on both games, consistent with faster convergence of the gradient signal. On Breakout, both algorithms require more steps to form well-separated clusters, reflecting the greater complexity of the game's state space.

\begin{figure}[H]
    \centering
    \begin{subfigure}[b]{0.48\linewidth}
        \includegraphics[width=\linewidth]{figures/tsne_temporal_dqn_pong.png}
        \caption{DQN / Pong: 500k-step intervals.}
    \end{subfigure}
    \hfill
    \begin{subfigure}[b]{0.48\linewidth}
        \includegraphics[width=\linewidth]{figures/tsne_temporal_ddqn_pong.png}
        \caption{DDQN / Pong: 500k-step intervals.}
    \end{subfigure}
    \vskip\baselineskip
    \begin{subfigure}[b]{0.48\linewidth}
        \includegraphics[width=\linewidth]{figures/tsne_temporal_dqn_breakout.png}
        \caption{DQN / Breakout: 500k-step intervals.}
    \end{subfigure}
    \hfill
    \begin{subfigure}[b]{0.48\linewidth}
        \includegraphics[width=\linewidth]{figures/tsne_temporal_ddqn_breakout.png}
        \caption{DDQN / Breakout: 500k-step intervals.}
    \end{subfigure}
    \caption{Temporal evolution of the representation space across all four agents. Representations start diffuse and tighten progressively over training. DDQN shows earlier structural emergence on both games.}
    \label{fig:temporal}
\end{figure}

\subsection{Saliency Maps}

Grad-CAM saliency maps computed over final checkpoints reveal consistent attention patterns across all four agents: the ball position attracts the highest saliency, with secondary attention on the agent's own paddle. Both algorithms have correctly identified the task-relevant objects.

The key difference lies in concentration. DQN saliency maps show more diffuse activation spread across wider screen regions. DDQN maps show tighter, more localised activation centred on the ball and paddle. On Breakout specifically, DDQN shows greater attention to bricks in the ball's immediate trajectory, indicating a more forward-looking attentional strategy (Figures~\ref{fig:saliency_pong} and~\ref{fig:saliency_breakout}).

\begin{figure}[H]
    \centering
    \includegraphics[width=0.7\linewidth]{figures/saliency_pong.png}
    \caption{Grad-CAM saliency maps on Pong. Both agents attend primarily to the ball; DDQN attention is more spatially concentrated.}
    \label{fig:saliency_pong}
\end{figure}

\begin{figure}[H]
    \centering
    \includegraphics[width=0.7\linewidth]{figures/saliency_breakout.png}
    \caption{Grad-CAM saliency maps on Breakout. DDQN shows more focused attention on the ball and the bricks directly in its path.}
    \label{fig:saliency_breakout}
\end{figure}

\subsection{Dead Neurons and Cosine Similarity}

The fraction of chronically inactive neurons, those firing in fewer than 5\% of game states, is consistently higher for DQN than DDQN across both games (Figure~\ref{fig:dead_neurons}). DQN's noisier gradient signal drives more neurons into the ReLU dead zone, from which gradient-based recovery is impossible, reducing the effective dimensionality of the representation layer. DDQN's more stable updates keep a larger fraction of neurons active, making fuller use of the available 512 dimensions. This is consistent with the capacity loss findings of Lyle et al.\ \cite{lyle2022understanding}.

\begin{figure}[H]
    \centering
    \includegraphics[width=0.7\linewidth]{figures/dead_neurons.png}
    \caption{Dead neuron fraction over training. DQN consistently has a higher proportion of inactive neurons than DDQN across both games.}
    \label{fig:dead_neurons}
\end{figure}

Cross-game cosine similarity between mean Pong and Breakout representations is non-trivial for both algorithms, consistent with the shared ball-and-paddle visual structure producing aligned representational directions (Figure~\ref{fig:cosine}). DDQN shows marginally higher cross-game similarity, suggesting its representations preserve more of the shared low-level visual structure rather than being perturbed by overestimation-driven gradient noise.

\begin{figure}[H]
    \centering
    \includegraphics[width=0.7\linewidth]{figures/cosine_similarity.png}
    \caption{Cross-game cosine similarity between mean Pong and Breakout representations. DDQN shows higher alignment, consistent with less noise in its learned features.}
    \label{fig:cosine}
\end{figure}

\subsection{Layer-wise CKA Similarity}

To localise where cross-game similarity breaks down within the network, CKA is computed between the Pong and Breakout agents at each layer depth: conv1, conv2, conv3, and fc\_repr (Figure~\ref{fig:cka}).

\begin{figure}[H]
    \centering
    \includegraphics[width=0.7\linewidth]{figures/layer_similarity_cka.png}
    \caption{Layer-wise CKA similarity between Pong and Breakout agents at each depth. Early convolutional layers maintain high cross-game similarity; similarity drops sharply at the representation layer.}
    \label{fig:cka}
\end{figure}

The pattern is consistent across both algorithms: CKA remains high at conv1 and conv2, begins to fall at conv3, and drops substantially at fc\_repr. This provides direct quantitative evidence for the representational hierarchy hypothesised from the cosine similarity and t-SNE results: early convolutional layers extract general visual features that are shared across games, while the representation layer encodes the game-specific strategic information needed for reward prediction. The transition from shared to specialised representations occurs progressively with depth, with the sharpest drop at the fully connected layer where the network moves from pixel-level feature extraction to state-value encoding.

This layer-wise profile also has practical implications. Because early conv layers are functionally similar across games, they are strong candidates for freezing during cross-game transfer or fine-tuning. The fc\_repr layer, by contrast, carries predominantly game-specific information and would need to adapt substantially in any cross-game transfer setting.

% ============================================================
\section{Discussion}
\label{sec:analysis}
% ============================================================

\subsection{Where Similarity Ends and Game-Specific Structure Begins}

The most important finding from the representation analysis is that shared visual structure does not translate into shared representational geometry at the level of the 512-dimensional representation layer. Despite Pong and Breakout having the same ball, paddle, and physics, agents playing different games develop clearly distinct internal encodings. The t-SNE projections show clean separation by game across both algorithms, and the cross-game cosine similarity, while non-trivial, is far from the high alignment one would expect if visual similarity were the dominant factor.

The CKA depth profile sharpens this picture considerably. High similarity at conv1 and conv2 confirms that early layers do respond to the shared visual content. The drop at conv3 and the sharper drop at fc\_repr mark the point at which game-specific strategic demands take over. The network is not simply learning the same features at every layer and then making different decisions at the output: it is actively diverging midway through the forward pass, at the boundary where low-level visual processing gives way to state-value encoding.

The explanation lies in what the representation layer is actually learning. The shared visual features -- ball movement, paddle position, deflection angles -- are present in both games, but the strategic context that determines which of those features are \textit{reward-predictive} differs entirely. In Pong, the opponent's paddle position matters enormously. In Breakout, the spatial distribution of remaining bricks drives Q-value differences. The representation layer learns to encode what is useful for predicting reward in each game, and because those useful features differ, the representations diverge even when the raw pixels overlap.

\subsection{Algorithm as a Modulator of Representation Quality}

DQN and DDQN were included in this study primarily to provide two independent observations of the same representational phenomenon. The fact that both algorithms show the same game-level separation confirms that the finding is not an artefact of one particular training procedure.

Within each game, however, DDQN consistently produces higher-quality representations: tighter t-SNE cluster geometry, fewer dead neurons, more focused saliency maps, and higher cross-game cosine similarity. These advantages are best explained by DDQN's cleaner gradient signal. DQN's overestimation bias inflates training targets in states where Q-value estimates are already unreliable, introducing noise that accumulates over millions of updates. DDQN's decoupled target computation removes this noise source, and the representation layer organises itself more consistently as a result.

Critically, these representational differences emerge even though DQN and DDQN achieve near-identical task performance (Table~\ref{tab:results}). This decoupling of performance and representation quality is itself an important finding: reward alone does not reveal whether an agent has built a clean internal model or a noisy one.

\subsection{Implications for Transfer Learning}

The CKA depth profile identified here has direct implications for cross-game transfer. Because conv1 and conv2 maintain high cross-game CKA similarity, they represent the most transferable components of the network: features learned on one game would remain relevant when applied to the other. Freezing these layers during cross-game fine-tuning would preserve their generalisation while allowing the deeper layers to specialise to the new game's strategic demands.

The fc\_repr layer, by contrast, shows substantial game-specific encoding. An agent pre-trained on Pong and fine-tuned on Breakout would need to substantially reorganise its representation layer, since the strategic demands differ, while its early convolutional features would carry over relatively well. The stronger cross-game CKA similarity under DDQN at deeper layers suggests DDQN retains marginally more transferable content at depth, which may give it a small advantage in fine-tuning settings.

\subsection{Limitations}

Several limitations should be considered when interpreting these results.

\textbf{Single random seed.} All experiments were run with seed 42. Deep RL results can be sensitive to initialisation, and running multiple seeds -- typically three to five in the literature \cite{henderson2018deep} -- would be needed to confirm that the observed representational differences are robust rather than specific to this initialisation. This is the most significant methodological limitation of the study.

\textbf{Reduced replay buffer.} The buffer capacity of 100,000 transitions is one-tenth of the 1M used in the original DQN paper \cite{mnih2015human}, a practical concession to memory constraints. Both algorithms trained with the same reduced buffer, so the reduction does not confound the algorithm comparison, but some effects might be more pronounced under the original conditions.

\textbf{Two games.} The study covers only Pong and Breakout. These were chosen for their structural similarity, which suits the research questions, but replication across a broader set of game pairs would strengthen confidence in the generalisability of the findings.

\textbf{Qualitative t-SNE assessment.} Cluster compactness in the t-SNE plots is assessed visually. Future work should complement this with quantitative measures such as the silhouette score or Davies-Bouldin index computed in the original 512-dimensional space, providing a more rigorous basis for the compactness claims.

% ============================================================
\section{Conclusion}
\label{sec:conclusion}
% ============================================================

This study has provided a controlled analysis of how structurally similar yet strategically distinct games shape the internal representations of deep RL agents, and how the training algorithm modulates the quality of those representations. The key findings are:

\begin{enumerate}
    \item \textbf{Shared visual structure does not produce shared representations.} Despite Pong and Breakout sharing a ball, paddle, and deflection physics, agents playing different games develop clearly distinct internal encodings. Strategic demands, not visual similarity, drive representational divergence at the 512-dimensional representation layer.
    \item \textbf{Representational similarity is spatially localised within the network.} CKA analysis shows high cross-game similarity at early convolutional layers, falling progressively with depth and dropping sharply at the representation layer. The boundary between shared visual processing and game-specific strategic encoding can be pinpointed to the transition from convolutional to fully connected computation.
    \item \textbf{Some cross-game similarity persists.} Non-trivial cosine similarity between Pong and Breakout representations reflects genuine shared visual content in the early convolutional layers, consistent with the CKA depth profile.
    \item \textbf{DDQN produces more structured representations within each game.} Tighter t-SNE cluster geometry, fewer dead neurons, more focused saliency, and higher cross-game cosine similarity all point to DDQN's cleaner gradient signal producing more organised internal models.
    \item \textbf{DQN overestimates Q-values systematically.} Mean max Q-values are substantially higher under DQN than DDQN, replicating the bias of van Hasselt et al.\ \cite{van2016deep} under reduced-buffer conditions.
    \item \textbf{Performance and representation quality are decoupled.} Both algorithms achieve near-identical task performance, confirming that reward alone is an insufficient indicator of whether an agent has built a clean internal model.
\end{enumerate}

The central finding is that the line between shared and game-specific representations in deep RL agents corresponds closely to the line between shared visual features and distinct strategic demands, and that this line can be localised to specific layers using CKA. The early convolutional layers generalise; the representation layer specialises. This has direct implications for transfer learning and motivates further study of cross-game transfer, layer freezing, and representation geometry across a broader range of game types and training algorithms.

% ============================================================
\newpage
\begin{thebibliography}{99}
\addcontentsline{toc}{section}{References}

\bibitem{mnih2015human}
Mnih V, Kavukcuoglu K, Silver D, Rusu AA, Veness J, Bellemare MG, et al.
Human-level control through deep reinforcement learning.
\textit{Nature}. 2015;518(7540):529--533.

\bibitem{van2016deep}
Van Hasselt H, Guez A, Silver D.
Deep reinforcement learning with double Q-learning.
In: \textit{Proceedings of the AAAI Conference on Artificial Intelligence}; 2016. vol.~30.

\bibitem{zahavy2016graying}
Zahavy T, Ben-Zrihem N, Mannor S.
Graying the black box: Understanding DQNs.
In: \textit{Proceedings of the 33rd International Conference on Machine Learning}; 2016. p.~1899--1908.

\bibitem{van2008visualizing}
Van der Maaten L, Hinton G.
Visualizing data using t-SNE.
\textit{Journal of Machine Learning Research}. 2008;9(86):2579--2605.

\bibitem{kornblith2019similarity}
Kornblith S, Norouzi M, Lee H, Hinton G.
Similarity of neural network representations revisited.
In: \textit{Proceedings of the 36th International Conference on Machine Learning}; 2019. p.~3519--3529.

\bibitem{bellemare2013arcade}
Bellemare MG, Naddaf Y, Veness J, Bowling M.
The Arcade Learning Environment: An evaluation platform for general agents.
\textit{Journal of Artificial Intelligence Research}. 2013;47:253--279.

\bibitem{selvaraju2017grad}
Selvaraju RR, Cogswell M, Das A, Vedantam R, Parikh D, Batra D.
Grad-CAM: Visual explanations from deep networks via gradient-based localization.
In: \textit{Proceedings of the IEEE International Conference on Computer Vision}; 2017. p.~618--626.

\bibitem{sutton2018reinforcement}
Sutton RS, Barto AG.
\textit{Reinforcement Learning: An Introduction}. 2nd~ed.
Cambridge, MA: MIT Press; 2018.

\bibitem{lyle2022understanding}
Lyle C, Rowland M, Dabney W.
Understanding and preventing capacity loss in reinforcement learning.
In: \textit{International Conference on Learning Representations}; 2022.

\bibitem{henderson2018deep}
Henderson P, Islam R, Bachman P, Pineau J, Precup D, Meger D.
Deep reinforcement learning that matters.
In: \textit{Proceedings of the AAAI Conference on Artificial Intelligence}; 2018. vol.~32.

\bibitem{saxe2013exact}
Saxe AM, McClelland JL, Ganguli S.
Exact solutions to the nonlinear dynamics of learning in deep linear neural networks.
In: \textit{International Conference on Learning Representations}; 2014.

\end{thebibliography}

\end{document}
